D'una manera molt simplificada, podem dir que la trompeta és un instrument musical de vent en què les diferents notes són produïdes aplicant aire per un extrem (que es considera tancat a causa de la presència dels llavis del músic) i que s'emeten per l'altre, considerat obert.

Les notes produïdes corresponen a determinats harmònics associats a les ones estacionàries que s'originen a l'instrument. La trompeta consta també de tres pistons que, quan es premen, augmenten de manera efectiva la longitud i canvien les notes emeses.

\begin{apartat}{1}
Si la longitud total del tub que representa la trompeta és $l_0 = 0,975\ \mathrm{m}$, indiqueu quina és la longitud d'ona i la freqüència dels tres primers modes de vibració estacionaris que es poden generar a la trompeta.
\end{apartat}

\begin{apartat}{1}
Quan el músic fa sonar l'instrument mentre prem el segon pistó, produeix la nota \textit{si} de la tercera octava, de freqüència $f = 247\ \mathrm{Hz}$. Sabent que aquesta nota correspon al segon mode de vibració permès a la cavitat de l'instrument, quina és ara la longitud efectiva de la cavitat? Quin és el recorregut extra $\Delta l$ que fa l'aire dins de la trompeta quan es prem aquest pistó?
\end{apartat}

\blocetiquetat{Dada:}{Velocitat del so en l'aire, $340\ \mathrm{m\,s^{-1}}$}
